\chapter{中断和设备驱动程序}
\label{CH:INTERRUPT}

\indextext{驱动程序}
是操作系统中管理特定设备的代码：
它配置设备硬件，
指示设备执行操作，
处理产生的中断，
并与使用该设备的进程交互。
驱动程序代码可能比较棘手，
因为驱动程序与设备并发执行，
而且通常与使用该设备的进程并发执行。
此外，驱动程序必须理解设备的硬件接口，
这可能很复杂且文档不完善。

需要操作系统关注的设备通常可以
配置为生成中断（interrupt），这是陷入（trap）的一种类型。
内核陷入处理代码识别到设备
产生了中断并调用驱动程序的中断处理程序；
在 xv6 中，这个分发发生在 {\tt devintr} \lineref{kernel/trap.c:/^devintr/} 中。

许多设备驱动程序在两个上下文中执行代码：
\indextext{上半部分（top half）}
在进程的内核线程中运行，
以及 \indextext{下半部分（bottom half）}
在中断时执行。上半部分
通过系统调用（如 {\tt read} 和 {\tt write}）被调用，
这些调用希望设备执行 I/O。
这段代码可能会要求硬件开始一个
操作（例如，要求磁盘读取一个块）；然后代码等待
操作完成。最终设备完成
操作并产生中断。驱动程序的中断处理程序，
作为下半部分，
确定什么操作已经完成，
如果需要则唤醒等待的进程，
并告诉硬件开始下一个操作（如果有的话）。

\section{代码：控制台输入}

控制台驱动程序 \fileref{kernel/console.c}
是驱动程序结构的简单示例。
控制台驱动程序通过连接到 RISC-V 的 \indextext{UART（通用异步收发传输器）}
串口硬件接收人类输入的字符。
控制台驱动程序一次累积一行输入，
处理特殊输入字符如
退格字符和控制键-u。用户进程，如 shell，使用
{\tt read} 系统调用从控制台获取输入行。
当你在 QEMU 中向 xv6 输入时，你的按键通过
QEMU 模拟的 UART 硬件传递给 xv6。

驱动程序与之通信的 UART 硬件是一个 16550
芯片~\cite{ns16550a}，由 QEMU 模拟。在真实的计算机上，16550
会管理连接到终端或其他计算机的 RS232 串行链路。
当运行 QEMU 时，它连接到你的键盘和显示器。

UART 硬件对软件来说表现为一组 \indextext{内存映射（memory-mapped）}
控制寄存器。也就是说，有一些物理地址
连接到 UART 设备，因此加载和存储
与设备硬件交互而不是与 RAM 交互。
UART 的内存映射地址从 0x10000000 开始，即 {\tt UART0}
\lineref{kernel/memlayout.h:/UART0.0x/}。
UART 有一些控制寄存器，每个都是
一个字节的宽度。它们相对于 {\tt UART0} 的偏移量在
\lineref{kernel/uart.c:/define.RHR/} 中定义。例如，
{\tt LSR} 寄存器包含指示输入
字符是否等待驱动程序读取的位。这些
字符（如果有）可以从
{\tt RHR} 寄存器读取。每次读取一个字符，UART 硬件
就从等待字符的内部 FIFO 中删除它，
并在 FIFO 为空时清除 {\tt LSR} 中的``就绪''位。
要传输数据，驱动程序向 {\tt THR} 寄存器
写入一个字节，这会导致 UART 将该字节附加到
UART 将在 RS232 串行链路上发送的字节 FIFO 中。
UART 发送和接收硬件在很大程度上是
相互独立的。

Xv6 的 {\tt main} 调用 {\tt consoleinit}
\lineref{kernel/console.c:/^consoleinit/} 来初始化 UART
硬件。这段代码配置 UART 在接收到
每个输入字节时产生
接收中断，
并在 UART 完成发送每个输出字节时产生
\indextext{发送完成（transmit complete）}
中断 \lineref{kernel/uart.c:/^uartinit/}。

xv6 shell 通过 {\tt init.c} 打开的文件描述符
\lineref{user/init.c:/open..console/} 从控制台读取。
对 {\tt read} 系统调用的调用通过内核到达 {\tt
  consoleread} \lineref{kernel/console.c:/^consoleread/}。{\tt
  consoleread} 等待输入到达（通过中断）并
缓冲在 {\tt cons.buf} 中，将输入复制到用户空间，
并（在一整行到达后）返回到用户进程。如果用户
还没有输入完整的一行，任何读取进程都会在
{\tt sleep} 调用中等待
\lineref{kernel/console.c:/sleep..cons/}
（第~\ref{CH:SLEEP} 章解释了 {\tt sleep} 的细节）。

当用户键入一个字符时，UART 硬件请求 RISC-V
产生中断，这会激活
xv6 的陷入处理程序。
陷入处理程序调用 {\tt devintr}
\lineref{kernel/trap.c:/^devintr/}，
它查看 RISC-V 的 {\tt scause} 寄存器以发现
中断来自外部设备。
然后它询问一个称为 PLIC 的硬件单元
\cite{riscv:priv}
告诉它是哪个设备产生的中断
\lineref{kernel/trap.c:/plic.claim/}。
如果是 UART，{\tt devintr} 调用 {\tt uartintr}。

{\tt uartintr}
\lineref{kernel/uart.c:/^uartintr/}
从 UART 硬件读取任何等待的输入字符
并将它们传递给 {\tt consoleintr}
\lineref{kernel/console.c:/^consoleintr/}；它不会
等待字符，因为将来的输入会产生新的中断。
{\tt consoleintr} 的工作是在
{\tt cons.buf} 中累积输入字符
直到一整行到达。
{\tt consoleintr} 特殊处理退格字符和一些其他字符。
当新行到达时，{\tt consoleintr} 唤醒
等待的 {\tt consoleread}（如果有的话）。

一旦被唤醒，{\tt consoleread} 将观察到 {\tt
  cons.buf} 中有一整行，将其复制到用户空间，
并（通过系统调用机制）返回到用户空间。

\section{代码：控制台输出}

在连接到控制台的文件描述符上执行 {\tt write} 系统调用
最终会到达
{\tt uartputc}
\lineref{kernel/uart.c:/^uartputc/}。
设备驱动程序维护一个输出缓冲区（{\tt uart\_tx\_buf}）
以便写入进程不必等待 UART 完成
发送；相反，{\tt uartputc} 将每个字符附加到缓冲区，
调用 {\tt uartstart} 启动设备传输（如果它还没有
开始），然后返回。{\tt uartputc}
唯一等待的情况是缓冲区已满。

每次 UART 完成发送一个字节时，它会产生一个中断。
{\tt uartintr} 调用 {\tt uartstart}，它检查设备
确实已经完成发送，并将下一个缓冲的输出字符交给设备。
因此，如果一个进程向控制台写入多个字节，
通常第一个字节将由 {\tt uartputc} 调用
{\tt uartstart} 发送，其余缓冲的字节将由
{\tt uartintr} 调用 {\tt uartstart} 在传输完成
中断到达时发送。

需要注意的一个通用模式是通过缓冲和中断将设备活动
与进程活动解耦。即使没有进程在等待读取，
控制台驱动程序也可以处理输入；
后续的读取会看到该输入。同样，进程可以发送输出
而不必等待设备。这种解耦可以通过允许进程
与设备 I/O 并发执行来提高性能，
并且在设备较慢（如 UART）或需要立即关注（如回显
键入的字符）时尤为重要。这个想法有时称为 \indextext{I/O
  并发（I/O concurrency）}。

\section{驱动程序中的并发}

你可能已经注意到在 {\tt consoleread}
和 {\tt consoleintr} 中调用了 {\tt acquire}。这些调用获取一个锁，
保护控制台驱动程序的数据结构免受并发访问。
这里有三个并发危险：两个在不同 CPU 上的进程可能
同时调用 {\tt consoleread}；
硬件可能请求 CPU 在 CPU 已经在 {\tt consoleread} 内部
执行时传递控制台（实际上是
UART）中断；
以及硬件可能在 {\tt consoleread} 执行时
在不同的 CPU 上传递控制台中断。
第~\ref{CH:LOCK} 章解释了如何使用锁
确保这些危险不会导致不正确的结果。

驱动程序中需要注意并发的另一种方式是，一个
进程可能正在等待设备的输入，但指示输入到达的中断
可能在不同的进程（或根本没有进程）运行时到达。
因此中断处理程序不允许考虑它们中断的进程或代码。
例如，中断处理程序不能安全地使用当前进程的页表调用 {\tt copyout}。中断处理程序通常
做相对较少的工作（例如，只是将输入数据复制到缓冲区），
并唤醒上半部分代码来完成其余工作。

\section{定时器中断}

Xv6 使用定时器中断来维持其对
当前时间的概念，并在计算密集型进程之间切换。定时器
中断来自连接到每个 RISC-V CPU 的时钟硬件。Xv6
为每个 CPU 的时钟硬件编程以定期中断 CPU。

{\tt start.c} 中的代码
\lineref{kernel/start.c:/^timerinit/} 设置一些控制位
允许监督模式访问定时器控制
寄存器，然后请求第一个定时器中断。
{\tt time} 控制寄存器包含一个
硬件以稳定速率递增的计数值；这作为
当前时间的概念。{\tt stimecmp} 寄存器
包含 CPU 将产生定时器
中断的时间；将 {\tt stimecmp} 设置为 {\tt time} 的当前值
加上 {\it x} 将安排一个 {\it x} 时间单位后的中断。
对于 {\tt qemu} 的 RISC-V
仿真，1000000 个时间单位约等于十分之一秒。

定时器中断通过 {\tt usertrap} 或 {\tt kerneltrap}
和 {\tt devintr} 到达，就像其他设备中断一样。
定时器中断到达时 {\tt scause} 的低位设置为
5；{\tt trap.c} 中的 {\tt devintr} 检测这种情况
并调用 {\tt clockintr}
\lineref{kernel/trap.c:/clockintr/}。
后者递增 {\tt ticks}，
允许内核跟踪
时间的流逝。递增只发生在一个 CPU 上，以避免多个 CPU 时时间过得更快。
{\tt clockintr} 唤醒在 {\tt pause}
系统调用中等待的任何进程，
并通过写入
{\tt stimecmp} 安排下一个定时器中断。

{\tt devintr} 为定时器中断返回 2
以向 {\tt kerneltrap}
或 {\tt usertrap} 指示它们应该调用 {\tt yield} 以便
CPU 可以在可运行进程之间多路复用。

内核代码可能被强制通过 {\tt yield} 进行上下文切换的定时器中断中断，
这是 {\tt usertrap} 中早期代码在启用中断之前小心保存 {\tt
  sepc} 等状态的部分原因。这些上下文切换也意味着
内核代码必须知道它可能在毫无预警的情况下从一个 CPU 移动到另一个 CPU。

\section{现实世界}

Xv6，像许多操作系统一样，允许在执行内核时产生中断甚至上下文切换（通过 {\tt yield}）。这样做的原因是在运行很长时间的复杂系统调用期间保持快速响应时间。然而，如上所述，在内核中允许中断是一些复杂性的来源；因此，少数操作系统只允许在执行用户代码时产生中断。

在典型计算机上支持所有设备的完整功能
是一项巨大的工作，因为有很多设备，设备有很多
功能，而且设备和驱动程序之间的协议可能很复杂
且文档不完善。在许多操作系统中，驱动程序占用的代码
比核心内核还多。

UART 驱动程序通过读取 UART 控制寄存器一次读取一个字节的数据；
这种模式称为 \indextext{程序控制 I/O（programmed I/O）}，
因为软件驱动数据移动。程序控制 I/O 很简单，但
速度太慢，无法在高数据速率下使用。需要以高速移动大量
数据的设备通常使用 \indextext{直接内存访问（DMA, direct memory access）}。
DMA 设备硬件直接将传入的数据写入 RAM，并从 RAM 读取
传出数据。现代磁盘和网络设备使用 DMA。DMA
设备的驱动程序会在 RAM 中准备数据，然后使用对控制寄存器的
单次写入告诉设备处理准备好的数据。

当设备需要在不可预测的时间关注，而且频率不太高时，
中断是有意义的。但中断具有很高的 CPU 开销。因此
高速设备，如网络和磁盘控制器，使用技巧来减少对中断的需求。
一个技巧是为整批传入或传出的请求产生单个
中断。另一个技巧是让驱动程序完全禁用中断，并定期检查设备是否需要关注。这种技术
称为 \indextext{轮询（polling）}。如果设备以高频率执行
操作，轮询是有意义的，但如果设备大部分时间是空闲的，
它会浪费 CPU 时间。一些驱动程序根据当前设备负载
在轮询和中断之间动态切换。

UART 驱动程序首先将传入数据复制到内核中的缓冲区，
然后复制到用户空间。这在低数据速率下是有意义的，但这种
双重复制可能会显著降低生成或消费数据非常快的设备的性能。
一些操作系统能够直接在用户空间缓冲区和设备硬件之间
移动数据，通常使用 DMA。

如第~\ref{CH:UNIX} 章所述，控制台对
应用程序来说表现为常规文件，应用程序使用 \lstinline{read} 和 \lstinline{write} 系统调用读取输入和写入输出。
应用程序可能想要控制无法通过标准文件系统调用表达的设备方面（例如，
在控制台驱动程序中启用/禁用行缓冲）。Unix
操作系统为此类情况提供 \lstinline{ioctl} 系统调用。

计算机的某些用途需要对外部事件的``实时''响应：
即保证在有限时间内发生的响应。
例如，在安全关键系统中，错过截止时间可能导致
灾难。Xv6 不适合实时场景。除了其他问题外，
xv6 的调度器在决定接下来运行哪个进程时没有考虑实时截止时间，
而且 xv6 有很长的禁用中断的内核代码路径，因此它可能无法快速响应中断。
实时操作系统不仅要解决这些问题，还必须以一种允许分析最坏情况响应时间的方式构建。

\section{练习}

\begin{enumerate}

\item 修改 {\tt uart.c} 以完全不使用中断。你可能需要
修改 {\tt console.c} 。

\item 为以太网卡编写驱动程序。

\end{enumerate}
